% =====================================================================
% Discussion
% =====================================================================
\section{Discussion}
\label{sec:discussion}

We organise the discussion around eight findings supported by the
data in Section~\ref{sec:results}.

\subsection{Waypoint adherence is a hard distinction, not a continuum}
\label{sec:disc-waypoint}

Linear, cubic, and quintic splines hit every waypoint to
floating-point precision on every variant in the experimental sweep,
while the B-spline approximation deviates by 0.28--0.92\,m on the
mobile-robot variants and 0.40\,rad on the robot arm. Increasing
waypoint density does not eliminate this deviation, because the
B-spline's clamped knot vector forces interpolation only at the
endpoints; the interior of the curve is fundamentally an
\emph{approximation} of the input control polygon. This is the
defining trade-off of the formulation, not a tuning parameter.

For applications in which the waypoints are strict constraints
(picking a part at an exact location), a B-spline approximation is
inappropriate regardless of how smooth its derivatives may be. For
applications in which the waypoints are soft guidelines (smoothing a
collision-avoidance buffer), the trade-off may be acceptable or even
desirable.

\subsection{Linear interpolation is the only method with $C^1$
violations}
\label{sec:disc-linear-c1}

The \texttt{max\_velocity\_disc} column quantifies $C^1$ violation
directly: linear interpolation produces $\sim 1$\,m/s velocity
discontinuities at every waypoint regardless of geometry, while the
smooth methods sit at or below $10^{-4}$\,m/s---essentially
numerical noise. The standard ``zero acceleration / zero jerk'' rows
of the linear results in \cref{tab:baseline-metrics} are therefore
misleading: they characterise the analytic derivative on the open
intervals between waypoints, not the distributional sense in which
the second derivative contains Dirac impulses.

A controller that finite-differences the commanded position, as is
standard in cascade-controlled servomechanisms, would observe
acceleration spikes proportional to its sampling rate at every
waypoint. The mismatch between the analytic and physical
interpretation of the linear-interpolation metrics motivates
reporting \texttt{max\_velocity\_disc} alongside the other metrics
in any future study.

\subsection{Adding waypoints can worsen smoothness}
\label{sec:disc-density}

The density sweep results
(Section~\ref{sec:results-density}) show that on the mixed geometry
under chord-length parameterisation, peak acceleration and RMS jerk
both increase monotonically with waypoint count for every spline
method. The mechanism is direct: chord-length parameterisation keeps
each segment's nominal speed at 1\,m/s, so denser waypoints reduce
$h_i$ proportionally. Because spline curvature scales like
$1/h_i^2$ and jerk like $1/h_i^3$, denser inputs force larger time
derivatives.

This finding contradicts the naive intuition that more waypoints
always means smoother motion. The correct interpretation is that
denser waypoints buy \emph{geometric fidelity}---the trajectory
converges towards the polyline as $n$ grows---at the price of
larger time-derivative magnitudes. Practitioners should choose
waypoint density based on which property matters more in their
application, not on the assumption that one trade-off dominates.

\subsection{Geometry dominates the smoothness budget}
\label{sec:disc-geometry}

Switching from a smooth sinusoid to a sharp 90$^{\circ}$ zigzag at
fixed $n=12$ approximately doubles peak acceleration and triples RMS
jerk for cubic and quintic splines---a much larger effect than any
of the other axes we tested. Geometry is therefore the dominant
determinant of how aggressively each method has to bend.

The B-spline absorbs this geometric character much more gracefully
than the interpolating splines: its sharp-vs.-gradual peak
acceleration ratio is only 1.18, compared to 1.78 for the cubic and
1.24 for the natural-BC quintic. This is a direct consequence of the
B-spline's approximation property---it is allowed to soften the
corner---and explains why B-splines are preferred in trajectory
planning for paths with intentionally sharp features such as
obstacle avoidance.

\subsection{Time spacing matters less than expected}
\label{sec:disc-spacing}

The textbook recommendation to use chord-length parameterisation
\cite[Section 24.1]{chapra2021} is reasonable but not as critical as
density or geometry. Switching to uniform spacing on the same
12-waypoint mixed geometry leaves peak acceleration and RMS jerk
within $\sim 30\%$ across all methods. The most visible effect is on
peak speed under linear interpolation, where uniform spacing forces
the robot to traverse long diagonals in the same time as short
straight segments.

For practical use, chord-length parameterisation remains the
default but the penalty of uniform spacing is modest enough that
implementers facing constraints (such as fixed-rate trajectory
buffers) need not feel obligated to enforce it.

\subsection{The quintic spline's apparent inferiority is a
boundary-condition artefact}
\label{sec:disc-quintic-bc}

The most surprising result in the baseline study is that the quintic
spline produced higher peak acceleration and RMS jerk than the cubic
on both datasets despite its higher continuity class
(\cref{tab:baseline-metrics}). The boundary-condition sensitivity
study (Section~\ref{sec:results-bc-sens}) explains this anomaly:
under natural BC the quintic must start and stop at rest in both
velocity and acceleration, but the rest of the spline is fitted to a
chord-length parameterisation that implies a 1\,m/s nominal speed;
reconciling these two demands within the first and last segment
forces a rapid transition with $\sim 4.5$\,m/s$^3$ jerk spikes at the
endpoints.

Replacing the natural BC with cubic-derived clamped BC---reading
$v_0$, $v_n$, $a_0$, $a_n$ off the cubic spline solution and
applying them as quintic boundary values---reduces peak acceleration
by 33\% and RMS jerk by 64\%, making the clamped quintic the
smoothest of the three interpolating methods on this dataset.

The practical implication is twofold. First, published comparisons
of quintic splines should explicitly state the boundary condition,
not just the method. Second, for applications where end velocities
and accelerations are known (a robot moving from one continuous
trajectory to another), the clamped quintic is the recommended
choice.

\subsection{Analytic derivatives match numerical differentiation}
\label{sec:disc-validation}

Centered finite-difference estimates of the dense-sampled trajectory
agree with the analytic derivatives to better than $10^{-3}$\,m/s$^2$
for acceleration and $10^{-3}$\,m/s$^3$ for jerk on the cubic,
quintic, and B-spline interpolants. This validates the textbook
analytic formulas of Section~\ref{sec:bg-cubic}--\ref{sec:bg-bspline}
as correct for the intended use.

The single exception---cubic spline jerk, with $5.1 \times 10^{-1}$
error---is not an implementation bug. The cubic spline's jerk is a
piecewise constant function with step discontinuities at every
internal knot. The five-point centered finite-difference scheme
cannot resolve those steps and instead smooths each into a ramp of
width $\sim 5h$. The error magnitude reflects the height of the
worst step in the dataset, not any error in the analytic formula.

\subsection{Computational cost is not a deciding factor}
\label{sec:disc-compute}

Even the most expensive method (the quintic, which solves a dense
$2(n+1) \times 2(n+1)$ linear system before evaluating) finishes in
under one millisecond at every density studied, including dense
evaluation at 1000 sample points. The full 28-row experiment runs
end-to-end in approximately 30\,ms.

For the problem sizes typical of robotic trajectory generation
(tens to low hundreds of waypoints), computation time is not a
deciding factor. The interesting trade-offs are between waypoint
adherence, smoothness, and geometric fidelity, not between methods
of different asymptotic cost.
